% Part: incompleteness
% Chapter: arithmetization-syntax
% Section: coding-formulas

\documentclass[../../../include/open-logic-section]{subfiles}

\begin{document}

\olfileid{inc}{art}{frm}
\olsection{ترميز \printtoken{P}{formula}}

قد ثبت التعريف العودي البدائي لـ$\fn{Term}(\OLId{latin-ordinary}{x})$.
ونتخذه الآن أساسًا لتعريف نظيرته في
ال!!{formula}s، وهي $\fn{Frm}(\OLId{latin-ordinary}{x})$،
تعريفًا عوديًا بدائيًا.

\begin{prop}
العلاقة $\fn{Atom}(\OLId{latin-ordinary}{x})$ عودية بدائية؛ ومعناها أن $\OLId{latin-ordinary}{x}$
عدد غودل ل!!{formula} ذرية.
\end{prop}

\begin{proof}
لتمييز $\OLId{latin-ordinary}{x}$ بوصفه عدد غودل ل!!{formula} ذرية،
يكفي ويلزم تحقق إحدى الحالات الآتية:
% تصحيح صياغة كلاسيكية (C284-01/أ): حُذف تكرار الفعل في صدر التفصيل.
\begin{enumerate}
\item ثمة $\OLId{latin-ordinary}{n}$ و$\OLId{latin-ordinary}{j}<\OLId{latin-ordinary}{x}$ و$\OLId{latin-ordinary}{z}<\OLId{latin-ordinary}{x}$ تحقق $\len{\OLId{latin-ordinary}{z}}=\OLId{latin-ordinary}{n}$، ولكل $\OLId{latin-ordinary}{i}<\OLId{latin-ordinary}{n}$ تكون
  $\fn{Term}((\OLId{latin-ordinary}{z})_\OLId{latin-ordinary}{i})$ صادقة، و$\OLId{latin-ordinary}{x}=$
\[
\Gn{\Obj P^\OLId{latin-ordinary}{n}_\OLId{latin-ordinary}{j}(} \concat \fn{flatten}(\OLId{latin-ordinary}{z}) \concat \Gn{)}.
\]
\item ثمة $\OLId{latin-ordinary}{z}_\OLMathNumeral{1},\OLId{latin-ordinary}{z}_\OLMathNumeral{2}<\OLId{latin-ordinary}{x}$ تصدق فيهما $\fn{Term}(\OLId{latin-ordinary}{z}_\OLMathNumeral{1})$ و$\fn{Term}(\OLId{latin-ordinary}{z}_\OLMathNumeral{2})$،
  و$\OLId{latin-ordinary}{x}=$
\[
\Gn{{\eq}(} \concat \OLId{latin-ordinary}{z}_\OLMathNumeral{1} \concat \Gn{,} \concat \OLId{latin-ordinary}{z}_\OLMathNumeral{2} \concat{\Gn{)}}.
\]
  \tagitem{prvFalse}{$\OLId{latin-ordinary}{x} = \Gn{\lfalse}$.}{}
  \tagitem{prvTrue}{$\OLId{latin-ordinary}{x} = \Gn{\ltrue}$.}{}
\end{enumerate}
\end{proof}

\begin{prop}
\ollabel{prop:frm-primrec}
العلاقة $\fn{Frm}(\OLId{latin-ordinary}{x})$ عودية بدائية، ومفادها أن $\OLId{latin-ordinary}{x}$
عدد غودل ل!!a{formula}.
\end{prop}

\begin{proof}
متتالية الرموز~$\OLId{latin-ordinary}{s}$ !!a{formula} إذا وفقط إذا
أمكنت متتالية تكوين $\OLId{latin-ordinary}{s}_\OLMathNumeral{0}$، \dots،
$\OLId{latin-ordinary}{s}_{\OLId{latin-ordinary}{k}-\OLMathNumeral{1}}=\OLId{latin-ordinary}{s}$ من !!{formula}، تسجل
بناء~$\OLId{latin-ordinary}{s}$ من !!{formula}s ذرية بقواعد
التكوين. ونختار المتتالية بلا تكرار،
فيكون كل $\OLId{latin-ordinary}{s}_\OLId{latin-ordinary}{i}$ صيغة جزئية من~$\OLId{latin-ordinary}{s}$،
وطولها لا يتجاوز $\OLId{latin-ordinary}{k}=\len{\OLId{latin-ordinary}{x}}$،
وعدد غودل لأي عنصر فيها لا يتجاوز~$\OLId{latin-ordinary}{x}$،
ويساويه في العنصر الأخير. فيحصر رمز
المتتالية، مثلًا، بالحد العودي البدائي
$\OLId{latin-ordinary}{p}_{\OLId{latin-ordinary}{k}-\OLMathNumeral{1}}^{\OLId{latin-ordinary}{k}(\OLId{latin-ordinary}{x}+\OLMathNumeral{1})}$. ومن ثم تكون
الكمّات التي نحتاج إليها على متتاليات
التكوين محدودة بحدود عودية بدائية.
% تصحيح صياغة كلاسيكية (C284-01/ب): أُزيل الرسم المعيب «المرّمزة» بإعادة العبارة.
\end{proof}

\begin{prob}
فصّل برهان \olref[inc][art][frm]{prop:frm-primrec} على مثال
البرهان الأول لـ\olref[inc][art][trm]{prop:term-primrec}.
\end{prob}

\begin{prop}
\ollabel{prop:freeocc-primrec}
العلاقة $\fn{FreeOcc}(\OLId{latin-ordinary}{x},\OLId{latin-ordinary}{z},\OLId{latin-ordinary}{i})$ عودية بدائية. وهي تعني
أن الرمز رقم $\OLId{latin-ordinary}{i}$ في الصيغة ذات عدد غودل~$\OLId{latin-ordinary}{x}$
ورود حر للمتغير ذي عدد غودل~$\OLId{latin-ordinary}{z}$.
\end{prop}

\begin{proof}
برهانه متروك للتمرين.
\end{proof}

\begin{prob}
برهن \olref[inc][art][frm]{prop:freeocc-primrec}. ولك الاستناد إلى أن كل مقطع من !!a{formula} يكون !!a{formula} لا بد أن يكون
!!{formula} جزئية منها.
\end{prob}

\begin{prop}
الخاصية $\fn{Sent}(\OLId{latin-ordinary}{x})$ عودية بدائية؛ ومعناها أن $\OLId{latin-ordinary}{x}$
عدد غودل ل!!a{sentence}.
\end{prop}

\begin{proof}
إنما ال!!{sentence} !!a{formula} خالية من الورودات الحرة
  لل!!{variable}s. ومن ثم يصدق $\fn{Sent}(\OLId{latin-ordinary}{x})$ إذا وفقط إذا
\begin{multline*}
\fn{Frm}(\OLId{latin-ordinary}{x}) \land {}\\
\bforall{\OLId{latin-ordinary}{i}<\len{\OLId{latin-ordinary}{x}}}{\bforall{\OLId{latin-ordinary}{z}<\OLId{latin-ordinary}{x}}{}}\\
(\bexists{\OLId{latin-ordinary}{j}<\OLId{latin-ordinary}{z}}{\OLId{latin-ordinary}{z}=\Gn{\Obj v_\OLId{latin-ordinary}{j}}} \lif \lnot\fn{FreeOcc}(\OLId{latin-ordinary}{x},\OLId{latin-ordinary}{z},\OLId{latin-ordinary}{i})).
\end{multline*}
% تصحيح مصدر (C284-02): أُضيف شرط كون x رمز صيغة قبل شرط انعدام المتغيرات الحرة.
\end{proof}

\end{document}
